%*******************************************************
% Abstract
%*******************************************************
%\renewcommand{\abstractname}{Abstract}
\pdfbookmark[1]{Abstract}{Abstract}
\begingroup
\let\clearpage\relax
\let\cleardoublepage\relax
\let\cleardoublepage\relax

\chapter*{Abstract}
In structural proof theory, natural deduction ($\ND$) combines introduction and elimination rules,
whereas Gentzen's $\LK$ is defined by dual introduction forms.
Through the Curry-Howard correspondence, the $\lambda$-calculus and $\lambda\mu\tilde{\mu}$-calculus
capture the computational essence of normalization in these respective frameworks.
However, alternative logical systems exist, such as Parigot's \enquote{Free Deduction} ($\FD$),
which relies exclusively on elimination rules.
While the term assignment system and semantics of such a calculus can be understood
in terms of its translation to either $\ND$ or $\LK$, we lack a standalone
computational interpretation for such a framework.

In this thesis, we present a modern reconstruction of Free Deduction,
titled $\piFD$, as a process calculus grounded in linear logic---akin
to Negri's \enquote{Uniform Calculus for Classical Linear Logic} (UCL).
By interpreting logical axioms as concurrent session channels and elimination
rules as the fulfillment of a future or promise,
we obtain a framework where the rule of
\textsc{cut} serves as a fork primitive, only facilitating interaction between
concurrent processes but not between endpoints of a channel.
This approach yields an alternative foundation for the design of session-typed disciplines that
diverges from traditional Curry-Howard interpretations of Girard's linear logic.

We formalize the metatheory of $\piFD$ within Rocq, providing mechanized
proofs for type safety and the Church-Rosser property.
These results establish a rigorous foundation for reasoning about concurrent
communication in elimination-based systems.
Ultimately, this work deepens the Curry-Howard correspondence for parallel and
communication-based calculi, offering new insights into the relationship
between logical elimination and session-based concurrency.

\endgroup

\vfill